% 第8回　店舗設計・空間構成（記入例・模範解答）
\session{8}{2}{11月18日（水）}{店舗設計・空間構成}{AIと空間記述}{yes}

\goals{
  \item 空間をゾーンと要素（素材、光、音、サイン）に分けて記述できる
  \item 言葉による空間記述を、画像生成のプロンプトに変換できる
  \item 生成画像を見て、コンセプトとの一致点とずれを言語化できる
  \item 生成画像の扱い（イメージであること、著作権と商標への注意）を理解する
}

\afillin{コンセプト文}{「見るだけ、大歓迎。」子どもも自分も、自分のペースで過ごせる“公園ラウンジ”のディーラー}

\work[90mm]{ワーク1}{空間の要件と、ゾーンごとの記述}
\begin{promptbox}[空間を言葉にする]
次のコンセプトとペルソナをもとに、店舗の空間を記述してください。入口から順に、来店客が通るゾーンごとに、広さの印象、床と壁の素材、光の質、音、配置されている家具や展示物、スタッフとの距離感を具体的に書いてください。コンセプト：（コンセプト文）　ペルソナ：（要約）
\end{promptbox}
\begin{wstabh}{colspec={Q[l,wd=16mm,m] X[l,m] X[l,m] X[l,m] X[l,m]},row{2-Z}={ht=15mm}}
  ゾーン & 広さ・素材 & 光・音 & 家具・展示物 & スタッフとの距離 \\
  受付 & \af{入口すぐの小さなカウンター。木の天板} & \af{外光が入る。静か} & \af{サインカード（見ているだけ／相談したい）を置く台} & \af{カードを渡すだけ。声かけはしない} \\
  展示 & \af{天井が高く開放的。床は芝生風のマット} & \af{自然光と間接照明。BGMは小さく} & \af{展示車3台。価格の目安と装備を書いたパネル} & \af{離れた位置で待機。カードを見て近づく} \\
  商談 & \af{半個室のブース。布張りのソファ} & \af{暖色の照明。仕切りで声が漏れにくい} & \af{タブレット、見積もりの見本} & \af{横並びに座り、一緒に画面を見る} \\
  待合 & \af{キッズスペースとカフェカウンターを隣接} & \af{明るい。子どもの声があっても気にならない広さ} & \af{遊具、絵本、ソファ、授乳室} & \af{見守り担当が1名} \\
  整備 & \af{工場との間はガラス張り} & \af{作業の音は抑える} & \af{作業の進み具合を表示する画面} & \af{完了時に呼び出し} \\
  & & & & \\
\end{wstabh}

\work[60mm]{ワーク2}{画像生成のプロンプトに変換し、条件を変えて生成する}
\begin{promptbox}[画像生成のプロンプトに変換する]
上の空間記述のうち「（ゾーン名）」の部分を、画像生成AI向けのプロンプトに変換してください。視点（入口から見た、俯瞰など）、時間帯、素材、光、人の有無を含め、英語と日本語の両方で出してください。
\end{promptbox}
\begin{wstabh}{colspec={Q[c,wd=6mm,m] Q[l,wd=16mm,m] X[2,l,m] X[2,l,m] X[2,l,m]},row{2-Z}={ht=14mm}}
  & ゾーン & 変えた条件（視点・時間帯など） & 表現したかったこと & うまくいかなかったこと \\
  1 & \af{展示} & \af{入口から見た視点、昼、人なし} & \af{芝生の上に車が置かれた公園のような開放感} & \af{屋外の公園になってしまい、店内に見えない} \\
  2 & \af{展示} & \af{「屋内」「天窓」を追加、親子連れあり} & \af{子ども連れでも居心地のよい雰囲気} & \af{車のロゴが実在メーカーに似てしまった（不使用）} \\
  3 & \af{待合} & \af{俯瞰、夕方、キッズスペースとカフェ} & \af{親が子どもを見ながら座れる配置} & \af{遊具が大きすぎ、カフェとの距離が不自然} \\
\end{wstabh}

\work{ワーク3}{画像を並べて言葉にする}
\atwoboxes{コンセプトと一致している点}{違和感・ずれ}{30mm}
{自然光と芝生の床で、ディーラーらしくない明るさが出ている。ソファから展示車とキッズスペースの両方が見える配置は、ペルソナAの要望に合う。}
{店内なのか屋外なのかがわかりにくい。展示車が多すぎて「売り場」の印象が強い。スタッフの姿がなく、相談できる店に見えない。}
\abox[空間記述の修正点]{20mm}{「屋内」「天窓のある高い天井」を明記する。展示車は3台に減らす。カウンターにスタッフが1人いる様子を入れ、「相談もできる」ことを表す。}

\work[90mm]{スケッチ}{平面のゾーニング（手描き可。入口・動線・各ゾーンを書き込む）}
\agridbox{85mm}{%
  \draw[ans,thick] (3,3) rectangle (160,80);
  \draw[ans,thick,fill=ans!6] (5,5) rectangle (55,45) node[midway]{キッズスペース\\（芝生ラウンジ）};
  \draw[ans,thick,fill=ans!6] (5,50) rectangle (55,78) node[midway]{カフェ\\カウンター・ソファ};
  \draw[ans,thick,fill=ans!6] (62,5) rectangle (100,20) node[midway]{受付\\（サインカード）};
  \draw[ans,thick,fill=ans!6] (62,26) rectangle (122,58) node[midway]{展示（車3台）\\セルフで乗り込み可};
  \draw[ans,thick,fill=ans!6] (128,26) rectangle (158,78) node[midway]{商談\\ブース\\（半個室）};
  \draw[ans,thick,fill=ans!6] (62,63) rectangle (122,78) node[midway]{ガラス越しの整備工場};
  \draw[ans,line width=1.2pt,->] (81,-2) -- (81,5);
  \node at (92,1.5) {入口};
  \draw[ans,dashed,thick,->] (81,20) -- (81,24) -- (60,35) -- (56,35);
  \draw[ans,dashed,thick,->] (100,22) -- (125,40);
  \node at (140,12) {- - → 主な動線};
}

\begin{caution}{生成画像を使うときの注意}
\achk 実在の店舗、ブランドのロゴ、既存の作品に似た画像は使わない
\achk 提案資料に載せる場合は「AIによる生成イメージ」と明記する
\achk 画像はあくまでイメージ。寸法や法規の裏づけは別に必要
\end{caution}

\begin{nexttime}
  \item 空間コンセプト図（ゾーニングを示した平面スケッチ。手描き可）を仕上げる
  \item 生成画像2〜3点と、それぞれで表現したかったこと・うまくいかなかったことのメモ（上のワーク2）
\end{nexttime}
\aimemoA{空間の言語化、画像生成プロンプトへの変換、画像の生成}
{「入口から順に、ゾーンごとに広さ・素材・光・音・家具・スタッフとの距離を」「展示ゾーンを画像生成用プロンプトに（英語と日本語）」}
{空間記述は表の形に自分たちで整理し直した。画像は3回生成し、実在メーカーのロゴに似た1枚は使わないことにした。}
{生成画像に実在の店舗やロゴに似た部分がないかをグループ全員で確認した}

\begin{point}
  \item ゾーンごとの記述が、素材・光・音・家具・スタッフとの距離まで具体的に書けているかを見る。「おしゃれな」「落ち着いた」だけの記述は、具体化を促す。
  \item 生成画像は完成図ではない。「うまくいかなかったこと」から記述を直せているか（ワーク3の修正点）を重視する。
  \item 実在メーカーのロゴや既存の店舗に似た画像を使っていないか、提案資料に「AIによる生成イメージ」と明記する予定かを確認する。
  \item スケッチは寸法の正確さより、入口・ゾーン・動線が読み取れることを求める。記入例の配置は一例。
\end{point}
